% =========================================================================
% PROVENANCE (section-level)
%   status      : CODE_FACT unless marked otherwise
%   produced-by : Recon_scripts/2_Reconstruction.py
%   commit      : dd520a5 (external-code submodule pin)
%   scope       : Proc16 + Prompt16 data, MC16rd; Run 1 + Run 2; e and mu
%   note        : every selection below was transcribed from the pinned
%                 source.  Physics motivation is marked TODO wherever the
%                 code does not document it -- do not invent one.
% =========================================================================
\section{Event Reconstruction}
\label{sec:reconstruction}

The signal-side $B$ candidate is reconstructed exclusively as
$B^{0} \to D^{-} \ell^{+}$ with $D^{-} \to K^{+} \pi^{-} \pi^{-}$, from four
charged final-state particles and no slow pion.
The tag-side $B$ is not reconstructed: it is represented by the rest of the
event (ROE), from which the tag-side kinematic quantities used in this
analysis are derived.
All selections in this section are applied in the steering script
\texttt{Recon\_scripts/2\_Reconstruction.py} at the pinned analysis commit,
and are transcribed literally in Table~\ref{tab:recon-cuts}.

\subsection{Final-state particle selection}
\label{sec:recon:fsp}

\paragraph{Track quality.}
All four charged final-state particles are required to satisfy a common
track-quality requirement: $|dz| < 2$~cm and $dr < 0.5$~cm with respect to the
interaction point, a polar angle inside the CDC acceptance, and at least one
hit in each of the CDC and the SVD.
The impact-parameter requirements select tracks compatible with originating
from the interaction region, and the hit requirements ensure that both the
momentum measurement and the vertex position are constrained by hits rather
than extrapolated.
This requirement is applied at list-filling time and therefore precedes every
selection below.

\paragraph{Hadron identification.}
Kaons are required to have $\texttt{kaonIDNN} > 0.9$ and $p > 0.3$~GeV$/c$;
pions $\texttt{pionIDNN} > 0.1$ and $p > 0.1$~GeV$/c$.
The asymmetry is large and deliberate: the kaon requirement is tight while the
pion requirement is close to no requirement at all.
Since the $D^{-} \to K^{+} \pi^{-} \pi^{-}$ final state contains one kaon and
two pions, and the dominant combinatorial background comes from pions
misidentified as kaons rather than the reverse, the kaon requirement carries
most of the rejection while the pion requirement is kept loose to preserve
efficiency on the two pions simultaneously.
\AnalysisTBD{the threshold values themselves are not documented in the source.
Record the efficiency and fake rate at $\texttt{kaonIDNN} > 0.9$ and
$\texttt{pionIDNN} > 0.1$, and whether either was scanned}

\paragraph{Lepton identification.}
Electrons and muons are both required to have a neural-network identifier
above $0.9$, and are then subjected to different momentum and acceptance
requirements: electrons $p > 0.3$~GeV$/c$, muons $p > 0.6$~GeV$/c$ together
with \texttt{inKLMAcceptance}.
The higher muon threshold and the KLM acceptance requirement go together:
muon identification relies on penetration to the KLM, so a muon must have
enough momentum to reach it and must point into its acceptance for the
identifier to be meaningful.
\AnalysisTBD{motivation for the specific values $0.9$, $0.3$~GeV$/c$ and
$0.6$~GeV$/c$ is not documented; record the working points and their
efficiencies}

The lepton momentum requirement is applied to the \emph{pre}-bremsstrahlung
momentum: the $p > 0.3$~GeV$/c$ cut is made on \texttt{e+:uncorrected} before
\texttt{correctBremsBelle} builds the corrected list.
An electron that would pass the threshold only after photon recovery is
therefore already removed.
\AnalysisTBD{confirm that applying the electron momentum requirement before
the bremsstrahlung correction is intended, and quantify the efficiency
difference against applying it afterwards}

\paragraph{Bremsstrahlung recovery.}
Electron candidates are corrected with \texttt{correctBremsBelle} using
photons from the \texttt{cdc} standard photon list with
$0.05 < E_{\gamma} < 2.5$~GeV and $\texttt{beamBackgroundSuppression} > 0.2$,
within a $0.15$~rad cone, allowing multiple photons per electron.
Recovering these photons matters here because both fit observables depend on
the lepton momentum: $\pdl$ directly, and $\mmiss$ through the signal-side
four-momentum.
An uncorrected electron migrates in both.

Both the corrected and the pre-correction (\texttt{BFbrems\_}) electron
quantities are written to the ntuple.
The pre-correction variables are the ones used for electron truth matching and
for the PID corrections (Sec.~\ref{sec:truth} and
Sec.~\ref{sec:corrections}), because the correction alters the momentum at
which the PID tables must be evaluated.

\paragraph{Track momentum scale and energy loss.}
Track momenta are corrected using the
\texttt{tracking\_release08\_20260119\_EnergyMomentumScale} global tag, with
separate payloads for data and simulation.
The \texttt{central} scale factor is applied to the nominal reconstruction;
the ten systematic variations are computed and aliased but are \emph{not}
written to the ntuple in the pinned version of the script, because the
\texttt{track\_variables} list is commented out of every variable collection.
They are therefore unavailable offline.
\AnalysisTBD{decide whether the momentum-scale variations need to be
re-enabled in the ntuple for the tracking systematic in Sec.~\ref{sec:syst}}

\paragraph{Event-level lepton requirement.}
Exactly one lepton is required per event,
\texttt{nParticlesInList(e-:goode) + nParticlesInList(mu-:goodmu) == 1}.
This is a requirement on the whole event, not on the candidate, and it has two
consequences.
The electron and muon channels become disjoint by construction, so they can be
treated as independent samples and combined without double counting.
It also removes events containing a second identified lepton, which suppresses
the double-semileptonic topologies in which the reconstructed lepton and the
$D$ come from different $B$ mesons --- part of the combinatorial background of
Sec.~\ref{sec:truth}.
\AnalysisTBD{quantify the signal efficiency loss of the one-lepton requirement
and its rejection of combinatorial background}

\subsection{$D^{-}$ candidate reconstruction}
\label{sec:recon:d}

$D^{\pm}$ candidates are formed from a kaon and two pions inside a wide
window, $1.78 < m(K\pi\pi) < 1.96$~GeV$/c^{2}$, which is loose enough that the
window does not shape the mass distribution before the vertex fit.
A tree fit is then performed with \texttt{updateAllDaughters} enabled and
without an IP constraint, since a $D^{\pm}$ has a measurable flight distance
and should not be constrained to the interaction point.
Candidates are required to have a reduced vertex $\chi^{2}$ below 13.
The fit also supplies \texttt{D\_A1FflightDistanceSig\_IP}, the flight-distance
significance with respect to the IP, which is one of the classifier inputs of
Sec.~\ref{sec:mva}: a real $D$ travels a measurable distance before decaying,
whereas a random $K\pi\pi$ combination points back to the interaction point.
\AnalysisTBD{motivation for the value $\texttt{vtxReChi2} < 13$ is not
documented; record its efficiency}
Because the tree fit updates the daughter momenta, the mass window is applied
a second time after the fit; the second window is the one that defines the
regions retained on tape.
That second window is not contiguous:
\begin{equation}
  m(K\pi\pi) \in
  \underbrace{[1.79, 1.82]}_{\text{left sideband}} \cup
  \underbrace{[1.855, 1.885]}_{\text{signal region}} \cup
  \underbrace{[1.92, 1.95]}_{\text{right sideband}}\ \mathrm{GeV}/c^{2}.
  \label{eq:dmass-windows}
\end{equation}
The intervals between the signal region and the sidebands are deliberately
not reconstructed.
The motivation is that signal tails and physics backgrounds populate the
region immediately adjacent to the signal peak, whereas the more distant
sidebands are dominated by fake-$D$ combinations; excluding the intermediate
region therefore yields a cleaner fake-$D$ control sample
(\texttt{ANALYSIS\_DECISION}).

This choice predates the decision to perform a simultaneous signal-region
plus sideband fit (Sec.~\ref{sec:fit}).
Under the current fit strategy a pure control region is no longer required,
since the fit models whatever is present in each channel.
\AnalysisTBD{potential future change: re-include the masked intervals
$[1.82,1.855]$ and $[1.885,1.92]$ in the sideband channel, or narrow the
sidebands towards the signal region.  This requires a new reconstruction
campaign because the masked events are not on tape.  Estimated contamination
of the masked region is below 5\%; source for that number to be supplied}

Throughout this note, a query of the form \texttt{D\_M < 1.85} in the
analysis code should be read as the left sideband $[1.79, 1.82]$, and
\texttt{D\_M > 1.9} as the right sideband $[1.92, 1.95]$: the intervening
values do not exist in the ntuples.

\subsection{$D^{*}$ veto}
\label{sec:recon:dstveto}

Candidates in which the $D^{-}$ originates from a $D^{*-} \to D^{-} \pi^{0}$
decay are identified by combining a $D^{\pm}$ from the mass signal region with
a slow $\pi^{0}$ from the \texttt{eff50\_May2020Fit} list.
Both photon daughters are additionally required to satisfy
$\texttt{beamBackgroundSuppression} > 0.5$ and
$\texttt{fakePhotonSuppression} > 0.1$, using the MC16rd weight files.
The $\pi^{0}$ from $D^{*} \to D \pi^{0}$ is soft, so its photons sit in the
energy range where beam background and misreconstructed clusters are most
abundant; without these requirements the veto would fire on random photon
pairs and remove genuine signal.

Because a $\pi^{0}$ can be built from more than one photon pair, at most one
$D^{*}$ candidate is retained per event, the one whose mass difference is
closest to the nominal $D^{*}$--$D$ splitting of $0.14065$~GeV$/c^{2}$.
The resulting $\Delta M$ is stored as an event-level quantity.
The number of slow $\pi^{0}$ candidates, \texttt{nPi0}, is stored as well.
The veto itself is applied offline as
$\Delta M < 0.135$ or $\Delta M > 0.145$~GeV$/c^{2}$
(\texttt{utilities.Dst\_veto\_cut}); the reconstruction stores the variable
without cutting on it, so the complementary region remains available as the
$B \to D^{*} \ell \nu$-enriched control region of Sec.~\ref{sec:validation}.
\AnalysisTBD{slow-$\pi^{0}$ efficiency corrections for the $D^{*}$ veto
region await an updated MC16rd table (expected summer 2026); until then the
$D^{*}$ veto region is not used quantitatively}

\subsection{$B^{0}$ candidate and rest of event}
\label{sec:recon:b}

$B^{0}$ candidates are formed as
\texttt{B0 =norad=> D- $\ell$+ ?addbrems} and vertex-fitted with an IP
constraint and \texttt{updateAllDaughters} disabled, so the $D$ daughter
retains the momenta from its own fit.
The reduced vertex $\chi^{2}$ is required to be below 14.
The distance between the $D$ and the lepton at the $B$ vertex is computed with
\texttt{calculateDistance}, giving \texttt{B0\_D\_l\_DisSig}: a genuine
$B \to D \ell$ decay has both daughters emerging from one point, while a
combinatorial pair does not.

Everything not used on the signal side forms the rest of the event, from which
every tag-side quantity in this analysis is derived.
Its composition is therefore as important as the signal-side selection, and it
is built in three stages under a single mask (\texttt{my\_mask}).
First a loose mask admits tracks with $dr < 10$~cm, $|dz| < 20$~cm, inside the
CDC acceptance and with $E < 5.5$~GeV, together with photons of
$0.05 < E_{\mathrm{cluster}} < 5.5$~GeV.
Second, $V^{0}$ candidates are used to update the ROE with a mass-constrained
tree fit, so that the daughters of a $K^{0}_{S}$ or $\Lambda$ are replaced by
the reconstructed parent rather than being counted as two unrelated tracks.
Third, the mask is tightened: tracks must satisfy
$\texttt{pValue} \geq 0.0005$ and a transverse-momentum-dependent
impact-parameter requirement that becomes looser at low $p_{T}$, where the
resolution is worse; photons must satisfy
$|t_{\mathrm{cluster}}| < 2\,\sigma_{t}$ and $|t_{\mathrm{cluster}}| < 200$~ns
together with $\texttt{beamBackgroundSuppression} > 0.5$ and
$\texttt{fakePhotonSuppression} > 0.1$.
Curl tracks are tagged separately with \texttt{tagCurlTracks} on a list of all
tracks.

Two observations on the tight mask that should be checked before the next
production:
\AnalysisTBD{the tight track requirement includes
\texttt{nCDCHits>=0}, which is satisfied by every track and therefore has no
effect.  Confirm whether a nonzero threshold was intended}
\AnalysisTBD{the tight track requirement is written as
\texttt{A and B and C and [$p_{T}<0.15$ and \ldots] or [\ldots] or \ldots}\@.
If \texttt{and} binds more tightly than \texttt{or} in the basf2 cut parser,
this applies the quality requirements \texttt{A}, \texttt{B} and \texttt{C}
only to the lowest $p_{T}$ branch, leaving tracks above $0.15$~GeV$/c$ with no
\texttt{pValue} or acceptance requirement.  Verify the parsed expression in
basf2 and add explicit brackets either way}

The tag-side vertex is obtained with \texttt{TagV} using the same mask, with a
tube constraint, the \texttt{standard\_PXD} track-finding type, no required
PXD hits, and a KFit reduced-$\chi^{2}$ requirement of 5.
Its quality enters the selection through \texttt{TagVReChi2IP}.

The ROE-based quantities entering the analysis are
$M_{\mathrm{bc}}^{\mathrm{ROE}}$, $\Delta E^{\mathrm{ROE}}$, the ROE charge,
the ROE extra energy, the ROE track multiplicity, and the missing-mass
squared $\mmiss$ (\texttt{B0\_recMissM2}), which is one of the two fit
observables.
The second fit observable, $\pdl$ (\texttt{p\_D\_l}), is \emph{not} produced
by the steering script; it is constructed offline as the sum of the two
centre-of-mass momenta already stored in the ntuple,
\texttt{p\_D\_l = D\_CMS\_p + ell\_CMS\_p}, for example in
\texttt{Notebooks/Plotting/1\_Plot\_BFBDT\_sigMC.ipynb}.
\AnalysisTBD{move the definition of \texttt{p\_D\_l} into shared library code
so that both fit observables have a single, citable producing step}

\paragraph{Tag-side requirements.}
Four loose requirements are applied to the ROE:
$M_{\mathrm{bc}}^{\mathrm{ROE}} > 4$~GeV$/c^{2}$,
$|\Delta E^{\mathrm{ROE}}| < 5$~GeV,
$|Q^{\mathrm{ROE}}| < 3$, and $0 < \texttt{TagVReChi2IP} < 100$.
The first two require that the remaining particles are at least broadly
consistent with a $B$ decay, the charge requirement removes events where the
ROE has lost or gained several tracks, and the last requires that the tag
vertex fit converged at all.
All four are deliberately loose compared with the offline requirements of
Eq.~\eqref{eq:offline-cut}, which tighten
$M_{\mathrm{bc}}^{\mathrm{ROE}}$ to above 5~GeV$/c^{2}$ and
$\Delta E^{\mathrm{ROE}}$ to $[-4, 1]$~GeV.
Keeping them loose on tape preserves the sidebands in these variables, which
are needed both as classifier inputs and to build the control regions of
Sec.~\ref{sec:validation}; the analysis context records that the
reconstruction-level selections on the $D$, $B$ and ROE kinematics were left
loose enough to retain training statistics for the classifier.

\subsection{Wrong-charge sample}
\label{sec:recon:wc}

For on-resonance ($4S$) samples only, the same chain is repeated with the
lepton charge reversed
(\texttt{allowChargeViolation=True}), producing the wrong-charge (WC)
control sample used in Sec.~\ref{sec:bbbar}.
Apart from the charge requirement, the WC reconstruction applies the same
vertex, ROE and tag-vertex selections as the nominal one.
The WC sample is not reconstructed for off-resonance data.

\subsection{Candidate multiplicity and best-candidate selection}
\label{sec:recon:bcs}

\AnalysisTBD{candidate multiplicity after the full selection, and the
best-candidate criterion actually used in the nominal chain.  The helper
\texttt{utilities.apply\_mva\_bcs()} exposes a \texttt{bcs} argument
defaulting to \texttt{'vtx'}, but the nominal working point and the resulting
multiplicity have not been established from a source artifact}

\subsection{Summary of selections}
\label{sec:recon:table}

Table~\ref{tab:recon-cuts} lists every selection applied in the pinned
steering script together with the literal code expression and its line
number, separating the requirement from its motivation.

\begin{table}[p]
  \centering
  \caption{Selections applied in \texttt{Recon\_scripts/2\_Reconstruction.py}
  at analysis commit \texttt{dd520a5}.  Line numbers refer to that file.
  Motivations marked TODO are not documented in the source and have not been
  invented here.}
  \label{tab:recon-cuts}
  \scriptsize
  \setlength{\tabcolsep}{3pt}
  \renewcommand{\arraystretch}{1.15}
  \begin{tabular}{>{\raggedright\arraybackslash}p{0.17\textwidth} >{\raggedright\arraybackslash}p{0.42\textwidth} >{\raggedright\arraybackslash}p{0.05\textwidth} >{\raggedright\arraybackslash}p{0.24\textwidth}}
    \hline
    Object & Code expression & Line & Motivation \\
    \hline
    All tracks &
      \texttt{abs(dz)<2 and dr<0.5 and thetaInCDCAcceptance and nCDCHits>0 and nSVDHits>0} &
      57 & Standard track quality \\
    $\pi^{\pm}$ &
      \texttt{pionIDNN > 0.1 and p > 0.1} & 58 & Kaon requirement carries the rejection; threshold TODO (Sec.~\ref{sec:recon:fsp}) \\
    $K^{\pm}$ &
      \texttt{kaonIDNN > 0.9 and p > 0.3} & 59 & TODO \\
    $e^{\pm}$ (pre-brems) &
      \texttt{electronIDNN > 0.9} & 61 & TODO \\
    $e^{\pm}$ (post-$p$ cut) &
      \texttt{p > 0.3} & 127 & TODO \\
    $\mu^{\pm}$ &
      \texttt{muonIDNN > 0.9}, then \texttt{inKLMAcceptance and p > 0.6} & 62, 129 & KLM acceptance required for muon ID \\
    Brems $\gamma$ &
      \texttt{0.05 < clusterE < 2.5 and beamBackgroundSuppression > 0.2} & 64 & Bremsstrahlung recovery input \\
    Event &
      \texttt{nParticlesInList(e-:goode) + nParticlesInList(mu-:goodmu) == 1} & 139 & Exactly one lepton; makes $e$/$\mu$ channels disjoint \\
    $D^{\pm}$ (pre-fit) &
      \texttt{1.78 < M < 1.96} & 148, 150 & Wide window before the tree fit \\
    $D^{\pm}$ vertex &
      \texttt{vtxReChi2 < 13} & 169 & Vertex quality \\
    $D^{\pm}$ (post-fit) &
      \texttt{1.79<M<1.82 or 1.92<M<1.95 or 1.855<M<1.885} & 149, 169 & Signal region plus clean fake-$D$ sidebands; see Eq.~\eqref{eq:dmass-windows} \\
    $\pi^{0}$ (slow) &
      \texttt{eff50\_May2020Fit} plus \texttt{beamBackgroundSuppression>0.5} and \texttt{fakePhotonSuppression>0.1} on both daughters & 178--182 & $D^{*}$ veto input \\
    $D^{*\pm}$ veto &
      best candidate by \texttt{abs(massDiff\_0 - 0.14065)}; $\Delta M$ stored, cut applied offline & 188--193 & Veto applied offline so the complement stays usable as a control region \\
    $B^{0}$ vertex &
      \texttt{vtxReChi2 < 14} & 231--232 & Vertex quality \\
    ROE mask (loose) &
      \texttt{dr<10 and abs(dz)<20 and thetaInCDCAcceptance and E<5.5}; $\gamma$: \texttt{0.05<clusterE<5.5} & 259--260 & Initial mask before $V^{0}$ update \\
    ROE mask (tight) &
      $p_{T}$-dependent \texttt{dr}/\texttt{dz} ellipses, \texttt{pValue>=0.0005}; $\gamma$: \texttt{clusterE>0.05}, \texttt{abs(clusterTiming)<2*clusterErrorTiming}, \texttt{<200}, \texttt{beamBackgroundSuppression>0.5}, \texttt{fakePhotonSuppression>0.1} & 261--269 & Beam-background and fake-photon rejection in the tag side \\
    Tag side &
      \texttt{4 < roeMbc and -5 < roeDeltae < 5 and abs(roeCharge) < 3 and 0 < TagVReChi2IP < 100} & 346--349 & Loose, to retain classifier training statistics and preserve sidebands; tightened offline (Sec.~\ref{sec:recon:offline}) \\
    \hline
  \end{tabular}
\end{table}

\subsection{Offline tightening}
\label{sec:recon:offline}

Two further selections are applied offline rather than during
reconstruction, and both are tighter than their reconstruction-level
counterparts.
The nominal offline requirement is \texttt{utilities.offline\_cut}:
\begin{equation}
  M_{\mathrm{bc}}^{\mathrm{ROE}} > 5,\quad
  -4 < \Delta E^{\mathrm{ROE}} < 1,\quad
  d_{r}(B^{0}) < 0.1,\quad
  p_{\ell} < 4 .
  \label{eq:offline-cut}
\end{equation}
The generic-$B\bar{B}$ weight tuning of Sec.~\ref{sec:bbbar} applies the same
requirements \emph{except} $p_{\ell} < 4$, which it omits
(\texttt{5\_BBbkg\_weights\_optuna\_minuit.py}, lines 809--814).
\AnalysisTBD{confirm whether the omission of $p_{\ell} < 4$ in the
generic-$B\bar{B}$ tuning is intentional.  If it is not, the tuned weights
were derived in a slightly wider region than the one they are applied in}

\subsection{Superseded and test-only code}
\label{sec:recon:superseded}

For completeness, the following reconstruction code exists in the analysis
repository but is not part of the nominal chain:
\texttt{Recon\_scripts/1\_Reconstruction\_test.py} (test harness) and
\texttt{Recon\_scripts/Old\_scripts/} (separate per-channel and per-charge
steering scripts, superseded by the single combined script).
In addition, the pinned nominal script hard-codes a single input file and an
output file name (lines 45--47); these are overridden by the grid submission,
so the checked-in values are a local test configuration rather than the
production one.
